\documentclass[11pt]{article}
\usepackage{graphicx}
\usepackage{amsmath}
\usepackage{booktabs}
\usepackage[margin=1in]{geometry}
\usepackage{setspace}
\usepackage{xcolor}
\usepackage{hyperref}
\hypersetup{colorlinks=true, linkcolor=blue, urlcolor=blue}

\title{License-Cap Policy Evaluation: Notes on Interpretation\\
\large What the Fiscal-Transfer Numbers Mean, and a Colorado Puzzle}
\author{Working notes -- entry game-marijuana project}
\date{August 2026}

\begin{document}
\maketitle

\begin{abstract}
\noindent This note walks through exactly what two summary numbers from the fiscal-transfer table mean,
since both are easy to misread. \emph{Share of border-county revenue} is not an independently measured
quantity, it is an algebraic transform of the regression coefficient itself, and when Montana's share
(48.0\%) exceeds Colorado and Washington's (42.8\%) despite Montana's far smaller dollar total, that is a
statement about relative composition within each regime's own border economy, not about which regime
transfers more in absolute terms. \emph{Transferred per catchment resident} is explained with its
governing assumptions made explicit, along with a genuine measurement issue discovered while writing this
note: Colorado+Washington's and Montana's catchment populations are not fully independent, they share two
Idaho counties. A closing discussion section takes up Colorado specifically: its border counties are less
likely to permit retail marijuana despite standing to gain fiscally from it, a pattern read here as local
political considerations outweighing fiscal incentives, without resolving why, which is left open for
future work.
\end{abstract}

\begin{table}[htbp]
\centering
\caption{Fiscal Transfer, Capped vs.\ Moratorium Regime (Coefficients Pooled 2022--2025, Applied to 2025 Revenue)}
\label{tab:magnitude1}
\small
\begin{tabular}{lcc}
\toprule
& Colorado + Washington & Montana \\
& (capped regime) & (moratorium regime) \\
\midrule
Coefficient used & $+0.559$ & $+0.654$ \\
& (pooled 2022--2025, $p=0.002$) & (pooled 2022--2025, $p=0.036$) \\
2025 border-county revenue & \$357.5M & \$48.3M \\
\textbf{Attributable to border effect} & \textbf{\$153.1M} & \textbf{\$23.2M} \\
Share of border-county revenue & 42.8\% & 48.0\% \\
Catchment population$^{a}$ & 512{,}162 & 198{,}016 \\
\textbf{Transferred per catchment resident} & \textbf{\$298.83} & \textbf{\$117.19} \\
\bottomrule
\end{tabular}
\begin{minipage}{0.95\textwidth}
\vspace{0.3em}
\footnotesize
\textit{Notes:} $^{a}$Catchment population is the union of non-legal-state counties within 70 miles of any
border county in that regime (20 counties across UT/WY/NE/ID for Colorado+Washington; 10 counties across
ID/ND/WY for Montana), not a per-county sum, to avoid double-counting non-legal counties near more than
one border county \emph{within} the same regime. See Section~\ref{sec:catchment-caveats} for a residual
overlap \emph{between} the two regimes' catchments. Attributable revenue $= \text{actual border revenue}
\times \text{premium}/(1+\text{premium})$, where premium $=100(e^{\beta}-1)$. This is an accounting
decomposition of an already-estimated coefficient, not an independent causal magnitude.
\end{minipage}
\end{table}


\section{How ``Share of Border-County Revenue'' Is Actually Measured}
\label{sec:share-measurement}

It helps to be precise about this, because it is not a separate empirical measurement, it is a direct
algebraic transformation of the regression coefficient already reported in the row above it.

The logic starts from the regression itself. For Colorado and Washington, the estimated coefficient is
$\beta=0.559$ on a binary border indicator in a regression of log revenue per dispensary, pooled over
2022--2025, the same four-year window used for Montana below. Interpreting that coefficient causally for
a moment (with the caveats of Section~\ref{sec:coefficient-meaning} in mind), it says a border dispensary
earns $100(e^{0.559}-1)=74.9\%$ more revenue than an otherwise identical interior dispensary in the same
state and year. Call this the \emph{premium}.

Now suppose the actual revenue a border dispensary earns, call it $R_{\text{actual}}$, equals a
counterfactual amount it would have earned without any border effect, $R_{\text{counterfactual}}$, scaled
up by the premium:
\[
R_{\text{actual}} = R_{\text{counterfactual}} \times (1 + \text{premium}).
\]
Rearranging, $R_{\text{counterfactual}} = R_{\text{actual}} / (1+\text{premium})$, and the dollar amount
attributable to the border effect is the difference:
\[
R_{\text{actual}} - R_{\text{counterfactual}} = R_{\text{actual}}\left(1 - \frac{1}{1+\text{premium}}\right)
= R_{\text{actual}} \times \frac{\text{premium}}{1+\text{premium}}.
\]
\emph{Share of border-county revenue} is exactly this last fraction, $\text{premium}/(1+\text{premium})$,
expressed as a percentage. For Colorado+Washington, premium $=74.9\%$, so share
$=0.749/1.749=42.8\%$. For Montana, premium $=92.3\%$ (from $\beta=0.654$), so share
$=0.923/1.923=48.0\%$.

The important thing to notice: this share depends \emph{only} on the coefficient $\beta$, not on how much
actual border revenue exists. It would be exactly the same 42.8\% whether Colorado and Washington's
border counties earned \$357.5M or \$3.575M. It is a restatement of the coefficient on a more intuitive
scale, cents attributable per dollar actually earned, not a second, independently verified fact about
the data.

\section{What the Coefficient Means, and What a Higher Share for Montana Actually Says}
\label{sec:coefficient-meaning}

The coefficient $\beta$ is a log-point gap between border and interior counties in the same state and
year, holding population fixed (and, for Colorado+Washington, holding the state fixed too, since the
regression pools both states with a state fixed effect). It answers: conditional on being in the same
state, the same year, and having the same population, how much more does a border county's dispensary (or
its per-capita cannabis market, for Montana) earn than an interior one? It is not a measure of how big a
state's cannabis market is, nor of how many potential customers live nearby; those are controlled for
separately elsewhere in the paper (population and, for Result 3, the captive non-legal population).

So when Montana's share of border-county revenue (48.0\%) exceeds Colorado and Washington's (42.8\%),
this is \textbf{not} a claim that Montana transfers more money in total, it obviously does not, $23.2M$
against $153.1M$. It is a claim about \emph{composition within each regime's own border economy}: of
every dollar a Montana border county earns in cannabis revenue, about 48 cents would not exist without the
proximity to a non-legal-state population, versus about 43 cents in Colorado and Washington. Montana's
border-county cannabis economy is, proportionally, slightly \emph{more} dependent on cross-border demand
than Colorado and Washington's is, even though Colorado and Washington capture a much larger absolute
transfer. A useful (if imperfect) analogy: a small shop very close to a state line can do the large
majority of its business with out-of-state customers even if a much bigger shop farther from any line
does more out-of-state business in raw dollars. Both statements can be true simultaneously, and Table~1
reports both because they answer different questions, one about relative composition, one about absolute
size, and the paper should be careful never to let a reader conflate the two.

One further honesty point: because share is a monotonic function of $\beta$ alone, the fact that
Montana's share exceeds Colorado and Washington's is arithmetically guaranteed the moment Montana's
coefficient (0.654) exceeds Colorado and Washington's (0.559), it is not independent confirmatory
evidence of anything beyond the coefficient comparison already discussed in Table~1 of the companion
synthesis note. It should be presented as a more interpretable restatement of that comparison, not as a
second, separate finding.

\section{Transferred per Catchment Resident: What It Is, and Real Caveats}
\label{sec:catchment-caveats}

This figure takes the attributable-dollar amount from Section~\ref{sec:share-measurement} and divides it
by \emph{catchment population}, the total number of people living in non-legal-state counties within 70
miles of any of that regime's own border counties (computed as a union across counties, not a per-county
sum, so a shared non-legal county near two border counties in the same regime is not counted twice). The
result, \$298.83 per catchment resident for Colorado+Washington versus \$117.19 for Montana, is meant to
answer whether the raw dollar gap is just a market-size artifact. It substantially is not: catchment
population is only 2.6 times larger for Colorado+Washington, while the dollar transfer is 6.60 times
larger ($\$153.1M$ vs.\ $\$23.2M$), so the per-resident gap of roughly 2.55 times cannot be explained by
market size alone.

This number should be handled carefully, for several real reasons, not just as a formality:

\begin{itemize}
\item \textbf{It assumes every catchment resident is an equally plausible customer.} It treats a person
living 5 miles from the border identically to one living 69 miles away, and does not adjust for income,
age, car ownership, or how normalized cannabis use is in that particular community. It is a rough
intensity measure, not a per-person willingness-to-cross estimate.
\item \textbf{It ignores price and product differences between states.} If Colorado's dispensaries charge
different prices than Montana's, or carry a different product mix, the dollar figure conflates
``more people crossing'' with ``each crossing person spending more,'' and this measure cannot separate the
two, consistent with the paper's existing acknowledgment that it has no price or quantity data.
\item \textbf{The 70-mile cutoff is the same fixed threshold used throughout the paper, not a
per-resident travel-cost estimate.} Real travel friction (highway access, winter road conditions) varies
within that radius and is not captured by straight-line distance.
\item \textbf{A genuine measurement issue, found while preparing this note: the two catchments are not
fully independent.} Idaho borders both Washington and Montana, and two Idaho counties (66{,}104
residents combined) fall within 70 miles of \emph{both} a Washington border county and a Montana border
county. This means about 13\% of Colorado+Washington's catchment and about 33\% of Montana's catchment is
the same underlying population, a resident of these two Idaho counties is, geographically, a plausible
customer for either regime's border dispensaries. This does not invalidate the comparison, both regimes'
border stores really could be drawing on this population, but it means the two \$298.83 and \$117.19
figures are not measuring two mutually exclusive customer pools, and this should be disclosed rather than
left implicit if this number is used in the paper.
\end{itemize}

\section{Discussion: Colorado, Where Fiscal Revenue Loses to Something Else}
\label{sec:colorado-puzzle}

The rest of this note, and the companion synthesis document, establish that Colorado's border counties
sit on a substantial, real fiscal opportunity: dispensaries there earn a large, statistically solid
premium (Section~\ref{sec:coefficient-meaning}), amounting to tens of millions of dollars a year in
revenue that would not exist without proximity to Utah, Wyoming, and Nebraska. A local government
deciding whether to permit retail marijuana under Amendment 64's opt-out authority has a straightforward
fiscal incentive to say yes, especially in a border county.

And yet Colorado's border counties are \emph{less} likely to have ever permitted an operating dispensary
than interior counties, 43\% versus 72\% (the political-economy result in the companion synthesis note).
The fiscal case for opting in is, if anything, strongest exactly where local governments have opted out
most. Put plainly: in Colorado's border region, whatever local governments are actually maximizing, it
does not appear to be fiscal revenue, or at least not fiscal revenue alone. Something else is outweighing
it, plausibly local political cost: retail marijuana may simply be unpopular enough in this part of the
state, a more rural, conservative stretch bordering Utah, Wyoming, and Nebraska, that the political cost
of allowing it exceeds the fiscal benefit a local government would otherwise capture. Call this
\emph{political revenue} in the loose sense that a local government may be responding to constituent
preference and electoral consequence rather than to the county budget.

This framing is offered as an interpretation of a real, documented pattern, not as a finding in itself.
\textbf{We deliberately do not test or resolve why fiscal considerations lose out here.} Doing so
properly would require data this project does not currently have: county-level voting patterns on
marijuana-related ballot measures, demographic or ideological measures for these specific border
counties, or the actual municipal ordinance records themselves rather than the has-any-dispensary proxy
used throughout. Any of these would let a future study distinguish genuine political-preference resistance
from other explanations (for instance, that some of these counties are simply too small for any operator
to have found retail marijuana commercially viable regardless of local policy, a possibility the
population-restricted check in the companion note does not fully rule out). We leave this open explicitly
as a direction for future research, rather than reaching for an explanation the current data cannot
support.

\end{document}
